\documentclass[spanish,aspectratio=169,xcolor={dvipsnames,table}]{beamer}
\usepackage[utf8]{inputenc}
\usepackage[T1]{fontenc}
\usepackage[spanish,es-tabla]{babel}
\usepackage[scaled=.96]{helvet}
\renewcommand{\familydefault}{\sfdefault}
\usepackage{booktabs}
\usepackage{tikz}
\usepackage{natbib}
\usepackage{hyperref}
\usetikzlibrary{arrows.meta,positioning,shapes.geometric,fit}

\definecolor{itescaRedDark}{HTML}{8F1117}
\definecolor{itescaRed}{HTML}{D70F18}
\definecolor{itescaCharcoal}{HTML}{141414}
\definecolor{itescaSilver}{HTML}{D9D9D9}
\definecolor{itescaPaper}{HTML}{F5F5F5}
\definecolor{itescaInk}{HTML}{181818}
\definecolor{itescaMuted}{HTML}{626262}

\usetheme{default}
\setbeamertemplate{navigation symbols}{}
\setbeamersize{text margin left=.65cm,text margin right=.65cm}
\setbeamercolor{normal text}{fg=itescaInk,bg=white}
\setbeamercolor{frametitle}{fg=white,bg=itescaCharcoal}
\setbeamercolor{block title}{fg=white,bg=itescaRedDark}
\setbeamercolor{block body}{fg=itescaInk,bg=itescaPaper}
\setbeamercolor{alerted text}{fg=itescaRedDark}
\setbeamerfont{frametitle}{size=\large,series=\bfseries}
\setbeamertemplate{itemize item}{\textcolor{itescaRedDark}{$\blacktriangleright$}}
\hypersetup{colorlinks=true,urlcolor=itescaRedDark,linkcolor=itescaRedDark,citecolor=itescaRedDark}
\newcommand{\itescalogo}{ITESCA/assets-itesca/web/logo-itesca-contorno.png}
\newcommand{\monograma}{ITESCA/assets-itesca/itesca-monograma.png}
\newcommand{\fuente}[1]{\par\vspace{.1cm}{\tiny\color{itescaMuted}Fuente: #1}}
\newcommand{\clave}[1]{\textcolor{itescaRedDark}{\textbf{#1}}}

\setbeamertemplate{frametitle}{%
  \nointerlineskip\begin{beamercolorbox}[wd=\textwidth,ht=.88cm,dp=.22cm,leftskip=.30cm,rightskip=.30cm]{frametitle}
  \insertframetitle\hfill\raisebox{-.07cm}{\includegraphics[height=.37cm]{\itescalogo}}\end{beamercolorbox}
  {\color{itescaRed}\rule{\textwidth}{1.3pt}}}
\setbeamertemplate{footline}{\leavevmode\hbox{%
  \begin{beamercolorbox}[wd=.31\paperwidth,ht=2.6ex,dp=1ex,leftskip=.6em]{author in head/foot}\color{white}\scriptsize Martin Jonathan de la Cruz\end{beamercolorbox}%
  \begin{beamercolorbox}[wd=.49\paperwidth,ht=2.6ex,dp=1ex,center]{title in head/foot}\color{white}\scriptsize Fundamentos de Gestión Administrativa\end{beamercolorbox}%
  \begin{beamercolorbox}[wd=.20\paperwidth,ht=2.6ex,dp=1ex,rightskip=.6em plus 1fil]{date in head/foot}\color{white}\scriptsize\hfill\insertframenumber/\inserttotalframenumber\end{beamercolorbox}}}
\setbeamercolor{author in head/foot}{bg=itescaCharcoal}
\setbeamercolor{title in head/foot}{bg=itescaRedDark}
\setbeamercolor{date in head/foot}{bg=itescaCharcoal}

\title{Organización, entorno y estructura}
\subtitle{Modelo de Kast y Rosenzweig · Actividad 3}
\author{Martin Jonathan de la Cruz}
\institute{Instituto Tecnológico Superior de Cajeme · Maestría en Gestión Administrativa}
\date{Agosto de 2026}

\begin{document}
\begin{frame}[plain]
\begin{tikzpicture}[remember picture,overlay]
\fill[itescaRedDark] (current page.south west) rectangle ([xshift=.59\paperwidth]current page.north west);
\fill[itescaCharcoal] ([xshift=.59\paperwidth]current page.south west) rectangle (current page.north east);
\node[opacity=.08] at ([xshift=2.8cm]current page.center) {\includegraphics[width=.58\paperwidth]{\monograma}};
\node[anchor=north west] at ([xshift=.55cm,yshift=-.45cm]current page.north west) {\includegraphics[width=3.9cm]{\itescalogo}};
\end{tikzpicture}
\vspace*{1.55cm}\begin{minipage}{.54\paperwidth}
{\color{white}\fontsize{23}{27}\selectfont\bfseries Organización, entorno y estructura\par}
\vspace{.25cm}{\color{itescaSilver}\large Modelo de Kast y Rosenzweig}\par
\vspace{.45cm}{\color{white}\normalsize Martin Jonathan de la Cruz}\par
{\color{white!80}\small Maestría en Gestión Administrativa}\end{minipage}\hfill
\begin{minipage}{.34\paperwidth}\raggedleft\includegraphics[width=.95\linewidth]{\itescalogo}\\[.4cm]
{\color{white}\small Fundamentos de Gestión Administrativa}\\{\color{white!75}\footnotesize Actividad 3 · Agosto de 2026}\end{minipage}
\end{frame}

\begin{frame}{Propósito: leer la organización como sistema}
\begin{block}{Pregunta guía}
¿Cómo explicar una organización sin reducirla al organigrama ni aislarla de las condiciones que la rodean?
\end{block}
\begin{columns}[T]\begin{column}{.48\textwidth}
\begin{itemize}\item Identificar los \clave{cinco subsistemas} interdependientes.
\item Distinguir el entorno general del entorno de tarea.
\item Comparar estructuras y elegir la que responda mejor a la contingencia.
\end{itemize}\end{column}\begin{column}{.47\textwidth}
\begin{alertblock}{Idea central} La estructura ordena el trabajo; el entorno exige adaptación; el modelo sistémico conecta ambas dimensiones.\end{alertblock}
\end{column}\end{columns}
\fuente{\citet{kastRosenzweig1985,kastRosenzweig1972}; \citet{bright2019}.}
\end{frame}

\begin{frame}{Kast y Rosenzweig: organización como sistema abierto}
\begin{center}\begin{tikzpicture}[node distance=.65cm,font=\small]
\node[draw,rounded corners,fill=itescaSilver!35,text width=3.1cm,minimum height=1cm,align=center] (in) {\textbf{Entradas}\\recursos, personas, información};
\node[draw=itescaRedDark,very thick,circle,fill=itescaRed!8,right=of in,minimum size=2.5cm,align=center] (org) {\textbf{Sistema}\\organización};
\node[draw,rounded corners,fill=itescaSilver!35,text width=3.1cm,minimum height=1cm,right=of org,align=center] (out) {\textbf{Salidas}\\bienes, servicios, resultados};
\draw[-{Latex[length=3mm]},very thick,itescaRedDark] (in)--(org);\draw[-{Latex[length=3mm]},very thick,itescaRedDark] (org)--(out);
\draw[-{Latex[length=3mm]},thick,bend left=38] (out) to node[above,font=\scriptsize]{retroalimentación} (org);
\end{tikzpicture}\end{center}
\begin{block}{Implicación directiva} Una decisión sobre tecnología, personas o estructura altera al sistema completo y su capacidad de responder al ambiente.\end{block}
\fuente{\citet{kastRosenzweig1972,kastRosenzweig1985}.}
\end{frame}

\begin{frame}{Los cinco subsistemas del modelo}
\small\begin{columns}[T]\begin{column}{.48\textwidth}
\begin{block}{1. Metas y valores} Define propósito, criterios de éxito y legitimidad.\end{block}
\begin{block}{2. Técnico} Integra conocimiento, procesos, equipos y métodos para transformar insumos.\end{block}
\begin{block}{3. Psicosocial} Reúne personas, motivación, liderazgo, relaciones y cultura.\end{block}
\end{column}\begin{column}{.48\textwidth}
\begin{block}{4. Estructural} Distribuye tareas, autoridad, coordinación y reglas.\end{block}
\begin{block}{5. Gerencial} Planea, integra, controla y vincula el sistema con el entorno.\end{block}
\begin{alertblock}{Clave de lectura} Ningún subsistema es suficiente por sí solo: el desempeño depende de su \clave{congruencia}.\end{alertblock}
\end{column}\end{columns}
\fuente{\citet{kastRosenzweig1985}; material de unidad: \emph{Kast y Rosenzweig UI}.}
\end{frame}

\begin{frame}{El entorno organizacional}
\footnotesize
\begin{columns}[T]\begin{column}{.49\textwidth}
\begin{block}{Entorno general (macro)}
Fuerzas que afectan indirectamente a muchas organizaciones:
\begin{itemize}\item económicas; \item tecnológicas; \item político-legales; \item socioculturales; \item demográficas y naturales.\end{itemize}
\end{block}\end{column}\begin{column}{.47\textwidth}
\begin{block}{Entorno de tarea (micro)}
Actores con influencia directa sobre resultados:
\begin{itemize}\item clientes o usuarios; \item proveedores; \item competidores; \item reguladores; \item aliados y comunidad.\end{itemize}
\end{block}\end{column}\end{columns}
\begin{alertblock}{Gestión estratégica} Vigilar, interpretar y responder a estas señales reduce incertidumbre y orienta decisiones.\end{alertblock}
\fuente{\citet{openstaxEnvironment2019,bright2019}.}
\end{frame}

\begin{frame}{De la incertidumbre a la respuesta}
\begin{center}\begin{tikzpicture}[node distance=.6cm]
\node[draw,rounded corners,fill=itescaSilver!35,minimum width=2.75cm,minimum height=.85cm] (a) {Señal del entorno};
\node[draw,rounded corners,fill=itescaRed!8,minimum width=2.75cm,minimum height=.85cm,right=of a] (b) {Interpretación directiva};
\node[draw,rounded corners,fill=itescaSilver!35,minimum width=2.75cm,minimum height=.85cm,right=of b] (c) {Ajuste organizacional};
\draw[-{Latex[length=3mm]},very thick,itescaRedDark] (a)--(b);\draw[-{Latex[length=3mm]},very thick,itescaRedDark] (b)--(c);
\end{tikzpicture}\end{center}
\begin{block}{Ejemplo} Una innovación tecnológica puede exigir capacitar personal (psicosocial), rediseñar procesos (técnico), redefinir responsabilidades (estructural) y revisar prioridades (metas).\end{block}
\fuente{Elaboración propia con base en \citet{kastRosenzweig1985,openstaxEnvironment2019}.}
\end{frame}

\begin{frame}{Diseño estructural: decisiones básicas}
\begin{columns}[T]\begin{column}{.49\textwidth}
\begin{itemize}\item \clave{Especialización}: división del trabajo.
\item \clave{Departamentalización}: agrupación por función, producto, territorio o cliente.
\item \clave{Jerarquía}: cadena de mando y tramo de control.
\end{itemize}\end{column}\begin{column}{.47\textwidth}
\begin{itemize}\item \clave{Centralización}: dónde se toman las decisiones.
\item \clave{Formalización}: reglas y procedimientos.
\item \clave{Coordinación}: mecanismos para integrar unidades.
\end{itemize}\end{column}\end{columns}
\begin{alertblock}{Criterio contingente} No existe una estructura universalmente superior: debe ajustarse a estrategia, tamaño, tecnología y dinamismo ambiental.\end{alertblock}
\fuente{\citet{openstaxStructures2019,bright2019}.}
\end{frame}

\begin{frame}{Tipos de estructura: funcional y divisional}
\small\begin{columns}[T]\begin{column}{.48\textwidth}
\begin{block}{Funcional}
Agrupa por especialidad: finanzas, operaciones, mercadotecnia, recursos humanos.
\begin{itemize}\item \textbf{Ventaja:} eficiencia técnica y economías de especialización.
\item \textbf{Riesgo:} silos y coordinación lenta entre áreas.
\end{itemize}\end{block}\end{column}\begin{column}{.48\textwidth}
\begin{block}{Divisional}
Agrupa por producto, mercado, cliente o región.
\begin{itemize}\item \textbf{Ventaja:} cercanía a resultados y capacidad de respuesta.
\item \textbf{Riesgo:} duplicación de recursos entre divisiones.
\end{itemize}\end{block}\end{column}\end{columns}
\fuente{\citet{openstaxStructures2019}.}
\end{frame}

\begin{frame}{Tipos de estructura: matricial, equipos y red}
\small\begin{columns}[T]\begin{column}{.31\textwidth}
\begin{block}{Matricial} Combina doble lógica —función y proyecto/producto—. Favorece integración, pero puede generar ambigüedad de autoridad.\end{block}\end{column}
\begin{column}{.31\textwidth}\begin{block}{Por equipos} Organiza trabajo interdisciplinario alrededor de procesos o problemas. Acelera colaboración y aprendizaje.\end{block}\end{column}
\begin{column}{.31\textwidth}\begin{block}{Red o virtual} Coordina capacidades internas y externas mediante alianzas y tecnología. Gana flexibilidad; exige confianza y control.\end{block}\end{column}\end{columns}
\begin{alertblock}{Elección} A mayor complejidad e interdependencia, mayor necesidad de coordinación lateral y mecanismos orgánicos.\end{alertblock}
\fuente{\citet{openstaxStructures2019,bright2019}.}
\end{frame}

\begin{frame}{Comparación ejecutiva: estructura según la contingencia}
\scriptsize\centering\begin{tabular}{p{2.4cm}p{2.8cm}p{2.8cm}p{2.8cm}}\toprule
\textbf{Estructura} & \textbf{Conviene cuando} & \textbf{Fortaleza} & \textbf{Atención directiva}\\\midrule
Funcional & operaciones estables y especializadas & eficiencia experta & coordinar entre áreas\\
Divisional & múltiples productos, regiones o mercados & foco en resultados & evitar duplicidad\\
Matricial & proyectos complejos e interdependientes & integración & aclarar autoridad\\
Equipos/red & ambiente cambiante e innovación & agilidad y aprendizaje & sostener confianza y comunicación\\\bottomrule\end{tabular}
\vspace{.25cm}\begin{block}{Síntesis} El diseño efectivo alinea metas, tecnología, personas, estructura y exigencias ambientales; esa es la utilidad práctica de pensar sistémicamente.\end{block}
\fuente{Síntesis propia a partir de \citet{kastRosenzweig1985,openstaxStructures2019}.}
\end{frame}

\begin{frame}{Conclusiones}
\begin{enumerate}\item Kast y Rosenzweig permiten describir la organización como un \clave{sistema abierto} de subsistemas interdependientes.
\item El entorno no es un fondo estático: sus fuerzas y actores condicionan información, recursos, riesgos y oportunidades.
\item La estructura es una decisión de diseño; debe ser congruente con la contingencia y facilitar coordinación, no solo control.
\item Una gestión responsable revisa continuamente la relación entre estrategia, entorno y capacidad organizacional.
\end{enumerate}
\vspace{.3cm}\begin{alertblock}{Declaración} Se utilizó IA generativa como apoyo de redacción y organización; las fuentes y la revisión final fueron verificadas por el autor.\end{alertblock}
\end{frame}

\begin{frame}{Referencias}
\footnotesize\bibliographystyle{apalike}\bibliography{fundamentos-de-gestion-administrativa}
\end{frame}
\end{document}
